\textbf{Background:} Heartbeat classifiers can look stronger in aggregate than they are on rare rhythms, particularly when patient-specific morphology and class imbalance are pronounced. We examined whether a Morlet CNN--Conformer could improve overall discrimination under subject-separated testing while making its class-specific failures explicit. \textbf{Methods:} Forty-five non-paced MIT-BIH records were divided at the record level: 38 were used for development and seven held-out records containing 15{,}573 beats were reserved for final evaluation. Three adjacent beats were represented as 32-scale Morlet scalograms. The CNN--Conformer was compared with attention-only, CNN--LSTM, and feature-engineered baselines on the same held-out split. \textbf{Results:} The CNN--Conformer achieved 0.60 accuracy (95\% CI 0.59--0.60), 0.26 macro-F1 (95\% CI 0.26--0.27), and 0.68 weighted-F1 (95\% CI 0.67--0.68), the highest values for these metrics among the evaluated architectures. The attention-only comparator reached 0.27, 0.15, and 0.38, respectively. The gain was not uniform across rhythms: Normal and ventricular ectopic beats had F1-scores of 0.74 and 0.57, whereas SVEB and Fusion had F1-scores of 0.00 and 0.01. Post-hoc threshold sweeps increased minority-class recall only at very low precision. \textbf{Conclusions:} The hybrid model improved aggregate performance within this controlled subject-separated benchmark, but rare-rhythm recognition remained the dominant limitation. The archived pipeline and class-wise diagnostics provide a reproducible basis for testing imbalance-aware methods and external validation.